\documentclass{article}
\usepackage{graphicx} % Required for inserting images
\usepackage{german}
\usepackage{amsmath}
\usepackage{amssymb}
\usepackage{amsthm}
\usepackage{url}


\newcommand{\R}{\mathbb{R}}% reele Zahlen
\newcommand{\C}{\mathbb{C}}% komplexe Zahlen
\newcommand{\Q}{\mathbb{Q}}% rationale Zahlen


\newtheorem{thm}{Satz}
\newtheorem{cor}[thm]{Korollar}
\newtheorem{lem}[thm]{Lemma}
\newtheorem{prop}[thm]{Proposition}


\title{Rekursion und Existenzaussagen}
\date{\today}
\author{Valerie Höhle\\ \url{valerie.hoehle@uni-bielefeld.de}}
\begin{document}

\maketitle
\section{Rekursionen}

Rekursionen sind ein zentrales Hilfsmittel der diskreten Mathematik. Die Grundidee besteht darin, neue Werte mithilfe bereits bekannter kleinerer Werte zu bestimmen. Ein wichtiges Beispiel hierfür sind die Binomialkoeffizienten.

Für \(n,k\in\mathbb{N}_0\) mit \(k\leq n\) werden die Binomialkoeffizienten durch

\[
\binom{n}{k}
=
\frac{n!}{k!(n-k)!}
\]

definiert. Sie geben die Anzahl der Möglichkeiten an, aus einer \(n\)-elementigen Menge genau \(k\) Elemente auszuwählen. Darüber hinaus gelten die Beziehungen

\[
\binom{n}{k}
=
\binom{n}{n-k}
\]

sowie

\[
\binom{n}{k}
=
\frac{n(n-1)\cdots(n-k+1)}{k!}.
\]

Um rekursive Formeln auch für allgemeinere Parameter formulieren zu können, erweitern wir die Definition der Binomialkoeffizienten auf komplexe Zahlen. Dazu setzen wir zunächst

\[
\binom{0}{0}=1,
\qquad
0!=1,
\qquad
n^{\underline 0}
=
n^{\overline 0}
=
1.
\]

Die Festlegung \(\binom{0}{0}=1\) ist dadurch motiviert, dass die leere Menge genau eine Teilmenge besitzt, nämlich sich selbst.

\bigskip

\textbf{Definition 1.1}

Für \(r\in\mathbb C\) und \(k\in\mathbb Z\) definieren wir

\[
\binom{r}{k}
=
\begin{cases}
\dfrac{r(r-1)\cdots(r-k+1)}{k!}, & k\geq 0,\\[1ex]
0, & k<0.
\end{cases}
\]

Die abfallende Fakultät wird durch

\[
r^{\underline{k}}
=
r(r-1)\cdots(r-k+1)
\]

definiert.

Diese Verallgemeinerung ermöglicht es, die bekannten Eigenschaften der Binomialkoeffizienten auf einen wesentlich größeren Definitionsbereich zu übertragen. Insbesondere lässt sich damit die Pascalsche Identität als Polynomidentität formulieren und beweisen.


\end{document}
